\clearpage
\section{Calcul matriciel}
\renewcommand{\titrecourant}{Semaine 6}

\begin{boussole}
\textbf{Objectif de la semaine.} Multiplier des matrices en sachant que l'ordre compte, calculer un déterminant et dire ce qu'il mesure, décider si une matrice est inversible, calculer son inverse, et résoudre un système linéaire en distinguant les trois situations : solution unique, infinité, aucune.\\
\textbf{À la fin, je sais :}
\begin{enumerate}
 \item vérifier la compatibilité des tailles puis calculer $AB$ ligne par colonne ;
 \item calculer $\det$ en $2\times2$ et en $3\times3$ et l'interpréter (aire, volume, coplanarité) ;
 \item calculer $A^{-1}$ ($2\times2$ par la formule, $3\times3$ par comatrice) et le vérifier par $AA^{-1}=I$ ;
 \item résoudre $AX=b$ quand $\det A\ne0$ et analyser le cas $\det A=0$.
\end{enumerate}
\textbf{Temps estimé} : 5 à 6 heures. L'annexe sur les vecteurs propres est un prolongement \horsprog\ (la diagonalisation n'est pas au programme du module) ; elle n'est pas évaluée.
\end{boussole}

\subsection{Test d'entrée (5 minutes)}
\begin{enumerate}
 \item Calculer $2\cdot5-3\cdot4$.
 \item Résoudre $\begin{cases}x+y=3\\x-y=1\end{cases}$.
 \item Quelle est l'aire du parallélogramme de côtés $(2,0)$ et $(1,3)$ ?
 \item Que vaut $\frac1a\cdot a$ ? Pour quel $a$ le nombre $\frac1a$ n'existe-t-il pas ?
 \item Produit scalaire de $(1,2)$ et $(3,4)$.
\end{enumerate}
\textit{Réponses : $-2$ ; $x=2$, $y=1$ ; $6$ ; $1$, $a=0$ ; $11$.}

\subsection{Cours}
\begin{definition}[Produit matriciel]
Si $A$ est de taille $m\times p$ et $B$ de taille $p\times n$, alors $AB$ est de taille $m\times n$ avec $(AB)_{ij}=\sum_{k=1}^pA_{ik}B_{kj}$. Le produit n'existe que si le nombre de colonnes de $A$ égale le nombre de lignes de $B$.
\end{definition}
\begin{propriete}[Règles]
$(AB)C=A(BC)$ ; $A(B+C)=AB+AC$ ; $AI=IA=A$ ; en général $AB\ne BA$ ; et $AB=0$ n'implique pas $A=0$ ou $B=0$.
\end{propriete}
\begin{definition}[Déterminant]
$\det\begin{pmatrix}a&b\\c&d\end{pmatrix}=ad-bc$. En $3\times3$, développement selon une ligne avec le signe en damier $\begin{smallmatrix}+&-&+\\-&+&-\\+&-&+\end{smallmatrix}$ : $\det A=a_{11}C_{11}+a_{12}C_{12}+a_{13}C_{13}$, où $C_{ij}=(-1)^{i+j}\times$ (déterminant $2\times2$ obtenu en retirant ligne $i$ et colonne $j$).
\end{definition}
\begin{propriete}[Interprétation et produit]
En $2\times2$, $|\det|$ est l'aire du parallélogramme construit sur les colonnes ; en $3\times3$, le volume du parallélépipède. Le déterminant est nul si et seulement si les colonnes sont linéairement dépendantes (coplanaires en dimension 3). De plus $\det(AB)=\det A\det B$.
\end{propriete}
\begin{propriete}[Inverse]
$A$ carrée est inversible si et seulement si $\det A\ne0$. En $2\times2$ : $A^{-1}=\frac1{ad-bc}\begin{pmatrix}d&-b\\-c&a\end{pmatrix}$. En général $A^{-1}=\frac1{\det A}\,{}^t\mathrm{Com}(A)$ où $\mathrm{Com}(A)$ est la matrice des cofacteurs. Et $(AB)^{-1}=B^{-1}A^{-1}$ (ordre inversé).
\end{propriete}
\begin{demonstration}[formule $2\times2$]
$\begin{pmatrix}a&b\\c&d\end{pmatrix}\begin{pmatrix}d&-b\\-c&a\end{pmatrix}=\begin{pmatrix}ad-bc&0\\0&ad-bc\end{pmatrix}=(ad-bc)I$ ; on divise par $ad-bc\ne0$. De même dans l'autre sens.
\end{demonstration}
\begin{propriete}[Système carré]
Soit $AX=b$ avec $A$ carrée. Si $\det A\ne0$ : solution unique $X=A^{-1}b$. Si $\det A=0$ : soit aucune solution, soit une infinité.
\end{propriete}
\begin{methode}
\textbf{Résoudre un système.} (1) Écrire $AX=b$. (2) Calculer $\det A$. (3) Si $\ne0$ : $X=A^{-1}b$, puis \emph{vérifier} en reportant. (4) Si $=0$ : élimination (combinaisons de lignes) pour trancher entre incompatible et infinité de solutions.
\end{methode}
\begin{piege}
(i) Multiplier terme à terme. (ii) Écrire $(AB)^{-1}=A^{-1}B^{-1}$. (iii) Oublier un signe dans le damier. (iv) Appliquer la formule de l'inverse sans tester $\det A\ne0$. (v) Conclure \og une seule solution \fg{} sans avoir calculé $\det A$.
\end{piege}

\subsection{Exemples travaillés et cas limites}
\begin{exemple}[non-commutativité]
$A=\begin{pmatrix}1&2\\0&1\end{pmatrix}$, $B=\begin{pmatrix}1&0\\3&1\end{pmatrix}$.\\
\emph{Solution.} $AB=\begin{pmatrix}7&2\\3&1\end{pmatrix}$, $BA=\begin{pmatrix}1&2\\3&7\end{pmatrix}$. Ces deux matrices sont différentes, bien que les deux produits existent.
\end{exemple}
\begin{cas}[produit nul, tailles incompatibles]
$\begin{pmatrix}1&0\\0&0\end{pmatrix}\begin{pmatrix}0&0\\0&1\end{pmatrix}=\begin{pmatrix}0&0\\0&0\end{pmatrix}$ alors qu'aucun facteur n'est nul : on ne peut pas \og simplifier \fg{} par une matrice. Par ailleurs, une matrice $2\times3$ multipliée par une matrice $2\times3$ n'a pas de produit (3 colonnes contre 2 lignes) ; mais $2\times3$ par $3\times2$ donne une matrice $2\times2$.
\end{cas}
\begin{exemple}[déterminant $2\times2$ et $3\times3$]
$\det\begin{pmatrix}3&1\\5&2\end{pmatrix}=6-5=1$. Pour $M=\begin{pmatrix}1&2&0\\0&1&3\\2&0&1\end{pmatrix}$, selon la première ligne : $1\cdot(1\cdot1-3\cdot0)-2\cdot(0\cdot1-3\cdot2)+0=1+12=13$.
\end{exemple}
\begin{cas}[déterminant nul]
$N=\begin{pmatrix}2&0&1\\1&3&2\\1&1&1\end{pmatrix}$ : $\det N=2(3-2)-0+1(1-3)=2-2=0$. Les colonnes sont coplanaires : $C_1=-C_2+2C_3$ (vérification : $-(0,3,1)+2(1,2,1)=(2,1,1)$). Le volume du parallélépipède est nul et $N$ n'est pas inversible.
\end{cas}
\begin{exemple}[inverse $2\times2$]
$\begin{pmatrix}2&1\\5&3\end{pmatrix}^{-1}=\frac1{6-5}\begin{pmatrix}3&-1\\-5&2\end{pmatrix}=\begin{pmatrix}3&-1\\-5&2\end{pmatrix}$. Contrôle : $\begin{pmatrix}2&1\\5&3\end{pmatrix}\begin{pmatrix}3&-1\\-5&2\end{pmatrix}=\begin{pmatrix}1&0\\0&1\end{pmatrix}$. $\checkmark$
\end{exemple}
\begin{exemple}[inverse $3\times3$]
Pour $M$ ci-dessus ($\det M=13$), le calcul des cofacteurs donne $13\,M^{-1}=\begin{pmatrix}1&-2&6\\6&1&-3\\-2&4&1\end{pmatrix}$. Contrôle sur la première ligne de $M$ : $(1,2,0)\cdot(1,6,-2)=13$ (colonne 1 de la matrice ci-dessus), et $(1,2,0)\cdot(-2,1,4)=0$. $\checkmark$
\end{exemple}
\begin{cas}[matrice singulière]
$\begin{pmatrix}2&4\\1&2\end{pmatrix}$ : $\det=4-4=0$. La seconde ligne est la moitié de la première. La formule de l'inverse divise par $0$ : pas d'inverse. Elle écrase le plan sur une droite (aire nulle).
\end{cas}
\begin{exemple}[système à solution unique]
$M\begin{pmatrix}x\\y\\z\end{pmatrix}=\begin{pmatrix}3\\4\\3\end{pmatrix}$. Comme $\det M=13\ne0$, $X=\frac1{13}\begin{pmatrix}1&-2&6\\6&1&-3\\-2&4&1\end{pmatrix}\begin{pmatrix}3\\4\\3\end{pmatrix}=\frac1{13}\begin{pmatrix}13\\13\\13\end{pmatrix}=\begin{pmatrix}1\\1\\1\end{pmatrix}$. Contrôle : $M(1,1,1)^T=(3,4,3)^T$. $\checkmark$
\end{exemple}
\begin{cas}[trois situations pour $\det A=0$]
Les systèmes $\begin{cases}x+y=2\\2x+2y=4\end{cases}$ et $\begin{cases}x+y=2\\2x+2y=5\end{cases}$ ont la même matrice de déterminant $0$. Le premier a une infinité de solutions ($y=2-x$) ; le second n'en a aucune (la seconde équation donnerait $x+y=\frac52$). Seul le second membre les distingue. Et $\begin{cases}x+y=2\\x-y=0\end{cases}$ ($\det=-2\ne0$) n'a que $(1,1)$.
\end{cas}

\subsection{Exercices}
\begin{exo}[A]\label{c6e1}
Pour $C=\begin{pmatrix}2&1\\0&3\end{pmatrix}$ et $D=\begin{pmatrix}1&4\\2&0\end{pmatrix}$, calculer $CD$ et $DC$.
\end{exo}
\begin{exo}[A]\label{c6e2}
Calculer le déterminant puis l'inverse de $\begin{pmatrix}3&1\\5&2\end{pmatrix}$.
\end{exo}
\begin{exo}[A]\label{c6e3}
Calculer $\det\begin{pmatrix}1&2&3\\0&1&4\\5&6&0\end{pmatrix}$.
\end{exo}
\begin{exo}[B]\label{c6e4}
Pour quelles valeurs de $m$ la matrice $\begin{pmatrix}1&m\\m&4\end{pmatrix}$ est-elle inversible ? Donner l'inverse en fonction de $m$.
\end{exo}
\begin{exo}[B]\label{c6e5}
Résoudre $\begin{pmatrix}2&1\\5&3\end{pmatrix}X=\begin{pmatrix}1\\2\end{pmatrix}$.
\end{exo}
\begin{exo}[B]\label{c6e6}
Un parallélogramme est construit sur $u=(3,1)$ et $v=(5,2)$. Calculer son aire. Que se passe-t-il pour $v=(6,2)$ ?
\end{exo}
\begin{exo}[C]\label{c6e7}
Discuter les solutions de $\begin{cases}x+y=2\\2x+2y=k\end{cases}$ selon $k$.
\end{exo}
\begin{exo}[C]\label{c6e8}
Avec $A$ et $B$ de l'exemple de non-commutativité, vérifier $\det(AB)=\det A\det B$ et $(AB)^{-1}=B^{-1}A^{-1}$.
\end{exo}
\begin{exo}[C]\label{c6e9}
Montrer que si $A$ est inversible et $AB=0$, alors $B=0$.
\end{exo}

\ifcorr
\subsection{Corrigés détaillés}
\begin{corrige}{\ref{c6e1}}
$CD=\begin{pmatrix}2\cdot1+1\cdot2&2\cdot4+0\\0+3\cdot2&0\end{pmatrix}=\begin{pmatrix}4&8\\6&0\end{pmatrix}$ ; $DC=\begin{pmatrix}2&1+12\\4&2\end{pmatrix}=\begin{pmatrix}2&13\\4&2\end{pmatrix}$. Donc $CD\ne DC$.
\end{corrige}
\begin{corrige}{\ref{c6e2}}
$\det=6-5=1$ ; $A^{-1}=\begin{pmatrix}2&-1\\-5&3\end{pmatrix}$. Contrôle : $\begin{pmatrix}3&1\\5&2\end{pmatrix}\begin{pmatrix}2&-1\\-5&3\end{pmatrix}=\begin{pmatrix}1&0\\0&1\end{pmatrix}$. $\checkmark$
\end{corrige}
\begin{corrige}{\ref{c6e3}}
Selon la première ligne : $1\cdot(0-24)-2\cdot(0-20)+3\cdot(0-5)=-24+40-15=1$.
\end{corrige}
\begin{corrige}{\ref{c6e4}}
$\det=4-m^2=(2-m)(2+m)$ : inversible si et seulement si $m\ne\pm2$. Alors $A^{-1}=\frac1{4-m^2}\begin{pmatrix}4&-m\\-m&1\end{pmatrix}$. Contrôle $m=0$ : $\frac14\begin{pmatrix}4&0\\0&1\end{pmatrix}=\mathrm{diag}(1,\frac14)$, inverse de $\mathrm{diag}(1,4)$. $\checkmark$ Pour $m=2$, les lignes $(1,2)$ et $(2,4)$ sont proportionnelles.
\end{corrige}
\begin{corrige}{\ref{c6e5}}
$\det=1$, inverse $\begin{pmatrix}3&-1\\-5&2\end{pmatrix}$ : $X=\begin{pmatrix}3-2\\-5+4\end{pmatrix}=\begin{pmatrix}1\\-1\end{pmatrix}$. Contrôle : $2-1=1$, $5-3=2$. $\checkmark$
\end{corrige}
\begin{corrige}{\ref{c6e6}}
Aire $=|\det\begin{pmatrix}3&5\\1&2\end{pmatrix}|=|6-5|=1$. Avec $v=(6,2)=2u$ : $\det=6-6=0$, les vecteurs sont colinéaires, le parallélogramme est aplati, d'aire nulle.
\end{corrige}
\begin{corrige}{\ref{c6e7}}
$\det=0$ pour tout $k$. La seconde équation s'écrit $x+y=\frac k2$ : si $k=4$, elle coïncide avec la première, infinité de solutions $(x,2-x)$ ; si $k\ne4$, incompatible avec $x+y=2$, aucune solution.
\end{corrige}
\begin{corrige}{\ref{c6e8}}
$AB=\begin{pmatrix}7&2\\3&1\end{pmatrix}$ : $\det=7-6=1=\det A\det B=1\cdot1$. $B^{-1}=\begin{pmatrix}1&0\\-3&1\end{pmatrix}$, $A^{-1}=\begin{pmatrix}1&-2\\0&1\end{pmatrix}$, $B^{-1}A^{-1}=\begin{pmatrix}1&-2\\-3&7\end{pmatrix}$, qui est l'inverse de $AB$ (formule $2\times2$ : $\begin{pmatrix}1&-2\\-3&7\end{pmatrix}$). $\checkmark$ Le produit $A^{-1}B^{-1}=\begin{pmatrix}7&-2\\-3&1\end{pmatrix}$ est différent.
\end{corrige}
\begin{corrige}{\ref{c6e9}}
Multiplions $AB=0$ à gauche par $A^{-1}$ : $A^{-1}AB=IB=B=A^{-1}\cdot0=0$.
\end{corrige}
\else\stepcounter{subsection}\fi

\subsection{Contrôle de fin de semaine (30 minutes, 20 points)}
\begin{enumerate}
 \item \bareme{6} Calculer $PQ$ et $QP$ pour $P=\begin{pmatrix}2&1\\0&3\end{pmatrix}$ et $Q=\begin{pmatrix}1&4\\2&0\end{pmatrix}$.
 \item \bareme{7} Montrer que $R=\begin{pmatrix}2&1\\5&3\end{pmatrix}$ est inversible et calculer $R^{-1}$.
 \item \bareme{7} Résoudre $RX=\begin{pmatrix}1\\2\end{pmatrix}$ et vérifier.
\end{enumerate}
\ifcorr
\textbf{Corrigé et barème.} \textbf{1.} $PQ=\begin{pmatrix}4&8\\6&0\end{pmatrix}$ \bareme{3}, $QP=\begin{pmatrix}2&13\\4&2\end{pmatrix}$ \bareme{3}. \textbf{2.} $\det R=1\ne0$ \bareme{3} ; $R^{-1}=\begin{pmatrix}3&-1\\-5&2\end{pmatrix}$ \bareme{4}. \textbf{3.} $X=R^{-1}b=(1,-1)^T$ \bareme{4} ; vérification $2-1=1$, $5-3=2$ \bareme{3}. Seuil conseillé : $14/20$.
\fi

\subsection*{Annexe \horsprog\ : valeurs et vecteurs propres}
Un vecteur $v\ne0$ est propre pour la valeur $\lambda$ si $Av=\lambda v$. Une valeur $\lambda$ est propre si et seulement si $\det(A-\lambda I)=0$ : cette condition assure l'existence d'un vecteur propre non nul, sans déterminer ce vecteur. Pour $S=\begin{pmatrix}2&1\\1&2\end{pmatrix}$ : $\det(S-\lambda I)=(2-\lambda)^2-1=(\lambda-1)(\lambda-3)$, donc $\lambda=1$ ou $\lambda=3$ ; vecteurs propres $(1,-1)$ pour $1$ et $(1,1)$ pour $3$. Contrôle : $S(1,1)^T=(3,3)^T$. Cas limite : $\begin{pmatrix}0&-1\\1&0\end{pmatrix}$ (rotation d'un quart de tour) n'a aucune valeur propre réelle ($\lambda^2+1=0$).

\subsection{Grille d'auto-évaluation}
$\square$ je vérifie les tailles avant de multiplier ; $\square$ je n'écris jamais $AB=BA$ sans preuve ; $\square$ je calcule $\det A$ avant de parler d'inverse ou d'unicité ; $\square$ je contrôle $AA^{-1}=I$ ; $\square$ pour $\det A=0$, je compare les seconds membres avant de conclure.
